\SourceLeaf{396}%
und
\[
\tfrac{c+d+e}{3}-\tfrac13\BareRoot{(cc+dd+ee+cd+ce+de)}
\]

Der Wurzelausdruck \(\BareRoot{(cc+dd+ee+cd+ce+de)}\) lässt sich auch in die Form
\[
\BareRoot{[(d-c)^2+(d-c)(e-d)+(e-d)^2]}
\]
bringen. Er entspricht folglich der dritten Seite eines Dreiecks mit den beiden Seiten \(d-c\) und \(e-d\) und dem eingeschlossenen Winkel \(=120^0\). Errichtet man also über \(CD\) ein gleichseitiges Dreieck mit der Spitze \(F\), dann ist \(EF\) gleich diesem Wurzelausdruck, sodass sich die beiden Werte von \(r\) leicht konstruieren lassen. Man kann leicht zeigen, dass der eine Wert zwischen \(c\) und \(d\), der andere zwischen \(d\) und \(e\) liegen muss und dass nur der erste tatsächlich dem Mittelpunkt der größten Ellipse entspricht. Für den anderen Wert erreicht nämlich
\[
\gamma+\delta r+\varepsilon rr+\zeta r^3
\]
kein Maximum, sondern ein Minimum oder vielmehr den größten \textit{negativen} Wert, dem also nur ein imaginärer Wert von \(\alpha\beta\) entsprechen kann. Man sieht leicht, dass sich dieser imaginäre Wert auf eine Hyperbel beziehen muss.

Sobald der Mittelpunkt der gesuchten Ellipse gefunden ist, lassen sich die übrigen unbekannten Größen ohne Schwierigkeit bestimmen. Aus \(\theta\) und \(r\) findet man \(t\). Aus \(t\) und \(u\) ergeben sich dann \(\alpha\) und \(\beta\) und aus einer oder mehreren der obigen Gleichungen \(\psi\). Damit sind sowohl die Abmessungen als auch die Lage der Ellipse vollständig bestimmt.

Ich muss noch darauf hinweisen, dass die hier gelöste Aufgabe von der kürzlich in der \textit{Monatl. Corresp.} gestellten etwas abweicht. Es gibt nämlich Fälle, in denen die größte Ellipse, die sich \textit{innerhalb} eines Vierecks zeichnen lässt, eine seiner vier Seiten nicht berührt. Diese Fälle näher zu betrachten, ist aber hier nicht meine Absicht.

\par\medskip\Needspace{4\baselineskip}
\pdfbookmark[2]{Direkter Beweis des obigen Satzes über Vierecke.}{ellipse-directer-beweis}
\begin{center}\textit{Direkter Beweis des obigen Satzes über Vierecke.}\end{center}
\hspace*{\parindent}Es seien \(A,B,C,D\) die vier Eckpunkte des Vierecks, \(E\) der Schnittpunkt von \(AB\) und \(DC\), \(F\) der Schnittpunkt von \(BC\) und \(AD\) und \(G,H\) und \(I\)
